\documentclass[11pt]{article}

% ============================================================
% INF3862 - ENTREGA 1
% Predicción de actividad física diaria
% Compilador: pdfLaTeX
% ============================================================

\usepackage[utf8]{inputenc}
\usepackage[T1]{fontenc}
\usepackage[spanish,es-noshorthands]{babel}
\usepackage{lmodern}
\usepackage{geometry}
\usepackage{microtype}
\usepackage{amsmath}
\usepackage{booktabs}
\usepackage{tabularx}
\usepackage{array}
\usepackage{longtable}
\usepackage{enumitem}
\usepackage{xcolor}
\usepackage{tikz}
\usepackage{float}
\usepackage{graphicx}
\usepackage{hyperref}
\usepackage{listings}
\usetikzlibrary{
    arrows.meta,
    positioning,
    shapes.geometric,
    calc
}

\geometry{left=1cm,right=1cm,top=1cm,bottom=1.6cm}

\definecolor{azul}{HTML}{17365D}
\definecolor{celeste}{HTML}{EAF1F8}
\definecolor{verde}{HTML}{28765C}
\definecolor{gris}{HTML}{555555}

\hypersetup{
    colorlinks=true,
    linkcolor=black,
    urlcolor=black
}

\setlength{\parindent}{1.2em}
\setlist[enumerate]{itemsep=1pt,topsep=2pt,leftmargin=*}
\setlist[itemize]{itemsep=1pt,topsep=2pt,leftmargin=*}

\newcommand{\seccion}[1]{\vspace{-1.2ex}\section*{#1}\vspace{-0.6ex}}

\newcommand{\codigo}[1]{\texttt{\small #1}}

\lstdefinestyle{sqlstyle}{
    basicstyle=\ttfamily\footnotesize,
    keywordstyle=\color{azul}\bfseries,
    commentstyle=\color{gris}\itshape,
    stringstyle=\color{verde},
    breaklines=true,
    showstringspaces=false,
    frame=single,
    rulecolor=\color{gray!35},
    backgroundcolor=\color{gray!4},
    columns=fullflexible,
    keepspaces=true,
    morekeywords={OVER,PARTITION,LEAD,ROWS,PRECEDING,COALESCE,LIMIT,WITH,BEGIN,COMMIT,ROLLBACK,RETURNING}
}
\lstset{style=sqlstyle}

\begin{document}

% ============================================================
% INFORME PRINCIPAL - PÁGINA 1
% ============================================================

\noindent
\begin{minipage}[c]{3.3cm}
  \includegraphics[height=3.95cm]{logo_uc.png}
\end{minipage}%
\begin{minipage}[c]{\dimexpr\textwidth-3.6cm\relax}
  Pontificia Universidad Católica de Chile\\
  Escuela de Ingeniería\\
  Magíster en Inteligencia Artificial\\
  INF3862 - Plataformas Big Data para Data Science\\
  Cuarto Bimestre 2026
\end{minipage}

\vspace{0.3cm}
\begin{center}
\textbf{Entrega 1: Diseño de datos}\\[2pt]
Predicción del nivel de actividad física diaria\\[2pt]
{\small Diego De la Vega \quad | \quad 16 de noviembre de 2026}
\end{center}

\seccion{1. Objetivo y preguntas guía}

Se diseña la plataforma de datos para \textbf{predecir los pasos
que dará cada usuario el día siguiente}, con la información recibida
antes del inicio de ese día. Si lo previsto queda bajo la meta del
usuario, la app le envía un recordatorio matinal con una actividad o
receta sugerida. No se buscan relaciones causales ni recomendaciones
clínicas.

{\small
\begin{tabularx}{\linewidth}{@{}p{2.1cm}X@{}}
\toprule
\textbf{Pregunta} & \textbf{Decisión del proyecto}\\
\midrule
Resultado &
Predicción diaria de pasos por usuario y dashboard histórico de
actividad, para el usuario de la aplicación y el equipo de producto.
Decisión: enviar o no el recordatorio matinal.\\
Datos &
Mediciones, sesiones y relojes del wearable; comidas registradas y
catálogo de recetas de Food.com. 61 días de recepción
(16-dic-2007 a 14-feb-2008), unos 470 MB en total (sección 2).\\
Operaciones &
Validar, normalizar los dos modelos de reloj, deduplicar, pasar a
fecha local, agregar por usuario-día, construir ventanas temporales
y predecir.\\
Exigencias &
Predicción lista antes de las 07:00 locales; lectura en la app en
menos de 2 s. Se tolera que un día quede incompleto por
sincronizaciones tardías, pero no duplicados en los agregados.
El modelo debe superar el error absoluto medio (MAE) del promedio
histórico del usuario.\\
Recursos &
Una persona y presupuesto acotado. Procesamiento local
(DuckDB/pandas sobre Parquet) antes de distribuir.\\
\bottomrule
\end{tabularx}
}

\textbf{Rol de Food.com.} Las comidas registradas y la energía de
las recetas preparadas (\codigo{Calories} por porción) entran como
variables de rutina diaria. Es una hipótesis a contrastar: si no
mejoran al modelo base, se registrará en el registro de decisiones.

\seccion{2. Fuentes, recepción y calidad}

{\small
\begin{tabularx}{\linewidth}{@{}>{\raggedright\arraybackslash}p{1.45cm}>{\raggedright\arraybackslash}p{3.1cm}>{\raggedright\arraybackslash}p{3.7cm}>{\raggedright\arraybackslash}X@{}}
\toprule
\textbf{Fuente} & \textbf{Registro y tipo} & \textbf{Volumen entregado} &
\textbf{Recepción (modo, frecuencia, latencia aceptable)}\\
\midrule
Recetas &
Una receta; estructurada, con listas &
522.517 filas; 179 MB Parquet &
Carga batch del catálogo, mensual; latencia de días.\\
Reseñas &
Una reseña; estructurada, con texto &
1.401.982 filas; 174 MB &
Batch; solo para validar \codigo{AuthorId} de los usuarios.\\
Relojes &
Un reloj y su vinculación; estructurado &
315 relojes, 300 usuarios; 20 KB &
Transacción en línea desde la app; menos de 1 s.\\
Mediciones &
Una lectura; JSONL semiestructurado &
15.540.132 filas; 1.449 archivos; 115 MB gz &
Lotes por sincronización, archivados por hora de recepción;
basta con que estén al corte diario.\\
Sesiones &
Una actividad; JSONL &
4.840 filas; 779 KB &
Evento al cerrar la sesión; minutos.\\
Comidas &
Una preparación; CSV estructurado &
22.106 filas; 1,0 MB &
Transacción en línea desde la app; menos de 1 s.\\
\bottomrule
\end{tabularx}
}

\textbf{Crecimiento.} Con 315 relojes llegan unas 255 mil lecturas
por día (cerca de 810 por reloj), es decir, 1,9 MB comprimidos al día
y unos 57 MB al mes. Para otra cantidad de relojes se escala con
$V_{\mathrm{dia}} = N_{\mathrm{relojes}} \cdot m_{\mathrm{dia}} \cdot b_{\mathrm{registro}}$,
con $b_{\mathrm{registro}} \approx 7{,}4$ bytes comprimidos por lectura.

\textbf{Problemas detectados en los datos entregados.}
\begin{itemize}
    \item \emph{Duplicados}: 345.599 mediciones repiten la clave
    (\codigo{device\_id}, \codigo{ts\_medicion}) (2,2\,\%), siempre con
    valores idénticos. Solo 45.584 son copias exactas de la fila; las
    300.015 restantes reaparecen en otro lote con distinto
    \codigo{ts\_recepcion}, por lo que se deduplica por esa clave. Sin deduplicar, los pasos
    se inflan 2,3\,\%. Hay además 48 sesiones duplicadas exactas.
    \item \emph{Atrasos}: entre la medición y la recepción hay una mediana
    de 1,2 h, un percentil 90 de 6,5 h, un percentil 99 de 32 h y un máximo
    de 105 h. El 12\,\% llega en un día posterior al de la medición.
    \item \emph{Formatos y faltantes}: FitBand 2 (13\,\% de las lecturas)
    usa otros nombres y fechas en epoch ms (Anexo~\ref{an:norm});
    frecuencia cardíaca nula en el 2,3\,\%; 48 comidas con
    \codigo{RecipeId} inexistente.
\end{itemize}

\textbf{Esquemas.} Los archivos crudos se guardan tal como llegan y
se leen con \emph{schema-on-read}, porque los dos modelos de reloj
difieren y conviene poder reprocesarlos. Las tablas operacionales y
analíticas usan \emph{schema-on-write}, para fijar tipos, claves y
reglas de calidad antes de consultar.


% ============================================================
% INFORME PRINCIPAL - PÁGINA 2
% ============================================================

\seccion{3. Modelo operacional (OLTP)}

{\small
\begin{tabularx}{\linewidth}{@{}>{\raggedright\arraybackslash}p{1.6cm}>{\raggedright\arraybackslash}X>{\raggedright\arraybackslash}p{4.4cm}@{}}
\toprule
\textbf{Entidad} & \textbf{Claves y columnas principales} &
\textbf{Modelo y motivo}\\
\midrule
Usuario &
\codigo{AuthorId} (PK), fecha de alta, meta de pasos. &
Relacional (FK).\\
Reloj &
\codigo{device\_id} (PK), modelo. &
Relacional.\\
Vinculación &
\codigo{vinculacion\_id} (PK), \codigo{device\_id} (FK),
\codigo{AuthorId} (FK), desde, hasta (nulo si activa), zona horaria. &
Relacional: la regla transaccional vive aquí.\\
Receta &
\codigo{RecipeId} (PK), nombre, categoría, \codigo{Calories},
proteína, porciones; ingredientes como lista JSON. &
Relacional con columnas documentales.\\
Comida &
\codigo{comida\_id} (PK), \codigo{AuthorId} (FK), \codigo{RecipeId} (FK),
fecha y hora de preparación, hora de registro, porciones $> 0$. &
Relacional: escritura puntual.\\
Sesión &
\codigo{session\_id} (PK), \codigo{device\_id} (FK), tipo, inicio, fin,
distancia, calorías. &
Relacional: lectura por reloj.\\
Medición &
(\codigo{device\_id}, \codigo{ts\_medicion}) (PK), recepción, pasos,
calorías, frecuencia cardíaca. &
Serie temporal \emph{append-only}: alto volumen, se lee por reloj y
rango de tiempo, fuera de la base de la app.\\
\bottomrule
\end{tabularx}
}

\textbf{Regla transaccional.} Un usuario tiene a lo más un reloj
activo, y un reloj está vinculado a lo más a un usuario. En los datos
hay 15 usuarios que cambiaron de reloj. El cambio exige cerrar la
vinculación anterior (\codigo{hasta = now()}) e insertar la nueva en
una misma transacción (consulta T1). Sin atomicidad, una falla entre
ambos pasos dejaría al usuario con dos relojes activos o sin ninguno:
sus pasos se contarían dos veces o se perderían. Dos índices únicos
parciales sobre \codigo{AuthorId} y \codigo{device\_id}
(\codigo{WHERE hasta IS NULL}) lo garantizan. El diagrama
entidad-relación está en el Anexo~\ref{an:er}.

\seccion{4. Modelo analítico (OLAP)}

{\small
\begin{tabularx}{\linewidth}{@{}p{3.6cm}X@{}}
\toprule
\textbf{Elemento} & \textbf{Definición}\\
\midrule
\codigo{fact\_actividad\_diaria} &
\textbf{Una fila por usuario y fecha local.} Claves:
\codigo{usuario\_key}, \codigo{fecha\_key}. Medidas: pasos totales,
calorías de actividad, minutos con lectura (cobertura), cantidad y
duración de sesiones.\\
\codigo{fact\_comida} &
\textbf{Una fila por comida registrada}, identificada por
\codigo{comida\_id} (la misma receta puede repetirse en un día). Claves:
\codigo{usuario\_key}, \codigo{fecha\_key}, \codigo{receta\_key}.
Medidas: porciones y energía (porciones $\times$ \codigo{Calories});
atributo: hora de registro.\\
\codigo{dim\_usuario} &
\codigo{AuthorId}, zona horaria, modelo del reloj activo, meta de pasos.\\
\codigo{dim\_fecha} &
Fecha, día de la semana, fin de semana, mes.\\
\codigo{dim\_receta} &
\codigo{RecipeId}, categoría, calorías, proteína y tiempo de preparación.\\
\bottomrule
\end{tabularx}
}

Las dos tablas de hechos comparten las dimensiones de usuario y fecha
(Anexo~\ref{an:dim}), así que la nutrición se agrega por usuario y día
sin multiplicar las filas de actividad. La energía de las recetas es
la \emph{preparada y registrada}, no el consumo confirmado. La
cobertura distingue un día de 0 pasos de un día sin el reloj puesto.

\textbf{Ubicación y formato.} Se usa un \emph{Lakehouse}: los datos
crudos (JSONL.gz, CSV y Parquet originales) quedan en almacenamiento
de objetos; las capas limpia y analítica, en Parquet columnar, y en
la E3 se evaluará un formato de tabla (Delta o Iceberg). El
almacenamiento columnar conviene porque las consultas analíticas leen
pocas columnas (pasos, fecha, usuario) de muchas filas.

\textbf{Partición.} Los crudos conservan la partición por fecha de
\emph{recepción}. La capa limpia se particiona por fecha local de
\emph{medición}, lo que favorece a A1 y A3, que filtran por rango de
fechas. Un lote tardío solo obliga a reescribir las particiones de
las fechas que toca. Como los pasos y las calorías de cada lectura son
incrementos desde la lectura anterior, el total diario es la suma tras
deduplicar. La fecha local se obtiene con la zona horaria de la
vinculación vigente en el momento de la medición.


% ============================================================
% INFORME PRINCIPAL - PÁGINA 3
% ============================================================

\seccion{5. Consultas y patrones de acceso}

El SQL, el usuario, la frecuencia, la latencia y la aceleración de
cada consulta están en el Anexo~\ref{an:sql}.

{\small
\begin{tabularx}{\linewidth}{@{}p{0.6cm}>{\raggedright\arraybackslash}p{1.9cm}>{\raggedright\arraybackslash}X>{\raggedright\arraybackslash}p{4.3cm}@{}}
\toprule
\textbf{ID} & \textbf{Clase} & \textbf{Propósito} & \textbf{Qué la acelera}\\
\midrule
T1 & OLTP, escribe &
Cambio de reloj: cierra la vinculación activa y crea la nueva
(\codigo{BEGIN ... COMMIT/ROLLBACK}). &
Índices únicos parciales.\\
T2 & OLTP, escribe &
Registrar una comida desde la app. &
PK e índices de FK.\\
T3 & OLTP, lee &
Reloj activo y zona horaria de un usuario. &
Índice por \codigo{(AuthorId, hasta)}.\\
T4 & OLTP, lee &
Últimas diez sesiones de un reloj. &
Índice por \codigo{(device\_id, inicio)}.\\
A1 & OLAP &
Promedio y total de pasos por usuario en un período (también es el
modelo de referencia). &
Parquet y poda por fecha.\\
A2 & OLAP &
Actividad diaria frente a energía de las recetas preparadas. &
Columnar; resumen usuario-día.\\
A3 & OLAP &
Dataset de predicción: ventanas de 7 días y pasos del día siguiente.
\textbf{Responde el objetivo.} &
Partición por fecha; variables materializadas.\\
\bottomrule
\end{tabularx}
}

Las consultas OLTP tocan una o pocas filas, son muy frecuentes y
concurrentes y deben responder en milisegundos. Las OLAP recorren
historiales completos con agregaciones y ventanas, corren en batch y
toleran segundos o minutos.

\textbf{Consulta fronteriza.} El promedio de pasos de los últimos 7
días, que la app muestra al abrirse, es analítico por su cálculo
(una ventana sobre el historial) pero operacional por su uso (miles de
lecturas concurrentes, menos de 2 s). Ejecutarlo sobre la base OLTP o
sobre las mediciones crudas competiría con la operación. Por eso se
calcula en el batch diario de la capa analítica, se publica en una
tabla resumen o caché indexada por usuario y la app solo la lee.

\seccion{6. Arquitectura inicial}

La aplicación escribe usuarios, vinculaciones y comidas en una base
operacional, que se extrae a diario hacia la zona cruda. Los archivos
de Food.com y los lotes de mediciones y sesiones llegan directo al
almacenamiento de objetos; el catálogo de recetas se carga además en
la base operacional, que lo necesita para validar comidas. Un proceso
diario normaliza los modelos de reloj, deduplica, pasa a fecha local y
valida referencias, primero hacia la capa limpia y luego hacia la
analítica. Esta alimenta el dashboard, el modelo y la tabla resumen
que lee la app. Los lotes tardíos reprocesan solo las fechas afectadas
mediante sobrescritura de partición, que es idempotente. El diagrama
está en el Anexo~\ref{an:arq}.

\seccion{7. Proyección}

En la E2 se implementará el pipeline sobre una muestra: reglas de
calidad de la sección 2 con su conteo, capa limpia en Parquet
particionado y consultas A1--A3. En la E3 se entrenará un modelo base
frente al promedio histórico, con separación temporal de 42, 9 y 9
días. Para evitar fugas, la predicción del día $d+1$ usará solo
lecturas recibidas y comidas registradas antes de las 00:00 locales;
el día $d$ llega incompleto al corte y así se reproducirá en
entrenamiento. Las variables usan un calendario continuo y se excluyen
los objetivos sin cobertura mínima, en vez de tratarlos como cero
pasos. Con 115 MB comprimidos, el caso no justifica hoy procesamiento
distribuido.

% ============================================================
% ANEXOS - NO CUENTAN EN LAS 3 PÁGINAS DEL INFORME
% ============================================================

\clearpage
\appendix

\begin{center}
{\Large\textbf{Anexos}}\\[3pt]
{\small Material de apoyo; no forma parte de las 3 páginas del informe.}
\end{center}
\vspace{-0.4cm}

\section{Modelo entidad-relación operacional}\label{an:er}

\begin{center}
\begin{tikzpicture}[
    entity/.style={
        draw=azul,
        fill=celeste,
        rounded corners,
        align=left,
        text width=4.1cm,
        inner sep=5pt,
        font=\footnotesize
    },
    rel/.style={-{Latex[length=2mm]},thick},
    lab/.style={font=\footnotesize, fill=white, inner sep=1pt}
]
\node[entity] (usuario) at (-5.6,0)
{\textbf{USUARIO}\\
PK: AuthorId\\
fecha\_alta, meta\_pasos};

\node[entity] (comida) at (0,0)
{\textbf{COMIDA}\\
PK: comida\_id\\
FK: AuthorId, RecipeId\\
fecha\_hora, registrado\_en\\
porciones};

\node[entity] (receta) at (5.6,0)
{\textbf{RECETA}\\
PK: RecipeId\\
nombre, categoría\\
Calories, proteína, porciones};

\node[entity] (vinc) at (-5.6,-3.4)
{\textbf{VINCULACIÓN}\\
PK: vinculacion\_id\\
FK: AuthorId, device\_id\\
desde, hasta (nulo = activa)\\
zona\_horaria};

\node[entity] (reloj) at (0,-3.4)
{\textbf{RELOJ}\\
PK: device\_id\\
modelo};

\node[entity] (sesion) at (5.6,-3.4)
{\textbf{SESIÓN}\\
PK: session\_id\\
FK: device\_id\\
tipo, inicio, fin\\
distancia\_km, calorías};

\node[entity] (medicion) at (0,-6.6)
{\textbf{MEDICIÓN}\\ \emph{serie temporal}\\
PK: device\_id, ts\_medicion\\
ts\_recepcion\\
pasos, calorías, FC};

\draw[rel] (usuario) -- node[lab,above]{1:N} (comida);
\draw[rel] (receta) -- node[lab,above]{1:N} (comida);
\draw[rel] (usuario) -- node[lab,right]{1:N} (vinc);
\draw[rel] (reloj) -- node[lab,above]{1:N} (vinc);
\draw[rel] (reloj) -- node[lab,above]{1:N} (sesion);
\draw[rel] (reloj) -- node[lab,right]{1:N} (medicion);
\end{tikzpicture}
\end{center}

La entidad MEDICIÓN es una estructura lógica. Por su volumen se
guarda en el almacenamiento de objetos y en la capa limpia, no en la
base transaccional de la aplicación.

\textbf{Reglas de integridad:}
\begin{itemize}
    \item Una comida referencia un usuario y una receta existentes
    (hoy fallan 48), con porciones mayores que cero.
    \item A lo más una vinculación activa por usuario y por reloj
    (índices únicos parciales; el cambio se hace con T1).
    \item Una sesión referencia un reloj válido y cumple
    \codigo{fin > inicio}; \codigo{session\_id} es único (hoy hay 48
    duplicados exactos).
    \item Una medición es única por (\codigo{device\_id},
    \codigo{ts\_medicion}) y cae dentro de un período de vinculación
    de su reloj.
\end{itemize}

\clearpage
\section{Arquitectura inicial}\label{an:arq}

\begin{center}
\begin{tikzpicture}[
    box/.style={
        draw=azul,
        fill=celeste,
        rounded corners,
        align=center,
        minimum height=1.05cm,
        text width=4cm,
        font=\footnotesize
    },
    src/.style={box, draw=verde, fill=green!7},
    arr/.style={-{Latex[length=2mm]},thick},
    lab/.style={font=\scriptsize, fill=white, inner sep=1pt}
]
% Fuentes
\node[src] (app) at (-5.6,0)
{\textbf{App móvil}\\Comidas y vinculaciones};
\node[src] (food) at (0,0)
{\textbf{Food.com}\\Recetas y reseñas (archivos)};
\node[src] (wear) at (5.6,0)
{\textbf{Relojes}\\Mediciones y sesiones};

% Recepción
\node[box] (oltp) at (-5.6,-2.3)
{\textbf{Base operacional}\\Transacciones en línea};
\node[box] (batch) at (0,-2.3)
{\textbf{Carga batch}\\Mensual};
\node[box] (lotes) at (5.6,-2.3)
{\textbf{Recepción de lotes}\\JSONL.gz por sincronización};

% Almacenamiento
\node[box] (raw) at (0,-4.6)
{\textbf{Almacenamiento de objetos}\\Crudo (bronze), por fecha de recepción};
\node[box] (silver) at (0,-6.6)
{\textbf{Normalización y calidad}\\Parquet limpio (silver), por fecha local};
\node[box] (gold) at (0,-8.6)
{\textbf{Modelo dimensional}\\Hechos, dimensiones y variables (gold)};

% Consumo
\node[box] (dash) at (-5.6,-10.9)
{\textbf{Dashboard}\\Actividad histórica};
\node[box] (ml) at (0,-10.9)
{\textbf{Modelo predictivo}\\Pasos del día siguiente};
\node[box] (cache) at (5.6,-10.9)
{\textbf{Tabla resumen / caché}\\Promedio 7 días y predicción};

\draw[arr] (app) -- (oltp);
\draw[arr] (food) -- (batch);
\draw[arr] (wear) -- (lotes);
\draw[arr] (batch) -- (raw);
\draw[arr] (lotes) |- (raw);
\draw[arr] (oltp) |- node[lab,pos=0.75,above]{extracción diaria} (raw);
\draw[arr,dashed] (batch) -- node[lab,above]{catálogo} (oltp);
\draw[arr] (raw) -- (silver);
\draw[arr] (silver) -- (gold);
\draw[arr] (gold) -| (dash);
\draw[arr] (gold) -- (ml);
\draw[arr] (gold) -| (cache);
\draw[arr] (ml) -- (cache);
\draw[arr,dashed] (cache.east) -- ++(0.45,0) |- node[lab,pos=0.25,right]{lectura de la app} ($(app.north)+(0,0.55)$) -- (app.north);
\end{tikzpicture}
\end{center}

\section{Modelo dimensional}\label{an:dim}

\begin{center}
\begin{tikzpicture}[
    dim/.style={
        draw=verde,
        fill=green!7,
        rounded corners,
        align=left,
        text width=3.9cm,
        inner sep=5pt,
        font=\footnotesize
    },
    fact/.style={
        draw=azul,
        fill=celeste,
        rounded corners,
        align=left,
        text width=5.3cm,
        inner sep=5pt,
        font=\footnotesize
    },
    arr/.style={-{Latex[length=2mm]},thick}
]
\node[fact] (f) at (-5.6,0)
{\textbf{FACT\_ACTIVIDAD\_DIARIA}\\
\emph{una fila por usuario y fecha local}\\
PK: usuario\_key, fecha\_key\\
pasos\_totales\\
calorias\_actividad\\
minutos\_con\_lectura\\
cantidad\_sesiones\\
duracion\_sesiones\_min};

\node[dim] (u) at (0,1.7)
{\textbf{DIM\_USUARIO}\\
usuario\_key, AuthorId\\
zona\_horaria\\ modelo\_reloj\\
meta\_pasos};

\node[dim] (d) at (0,-1.7)
{\textbf{DIM\_FECHA}\\
fecha\_key, fecha\\
dia\_semana, fin\_de\_semana\\
mes};

\node[fact] (c) at (5.6,0)
{\textbf{FACT\_COMIDA}\\
\emph{una fila por comida registrada}\\
comida\_id\\
usuario\_key, fecha\_key\\
receta\_key\\
porciones, energia\_kcal\\
registrado\_en};

\node[dim] (r) at (5.6,-3.6)
{\textbf{DIM\_RECETA}\\
receta\_key, RecipeId\\
categoría, Calories\\
proteína, tiempo\_total};

\draw[arr] (f) -- (u);
\draw[arr] (f) -- (d);
\draw[arr] (c) -- (u);
\draw[arr] (c) -- (d);
\draw[arr] (c) -- (r);
\end{tikzpicture}
\end{center}

Las dos tablas de hechos comparten dimensiones conformadas de usuario
y fecha. Así se cruza la nutrición con la actividad sin duplicar los
pasos de un mismo día.

\section{Consultas SQL y justificación}\label{an:sql}

\subsection{T1. Cambiar de reloj (escritura transaccional)}

\textbf{Usuario:} la aplicación, cuando el usuario vincula un reloj
nuevo. \textbf{Propósito:} reemplazar el reloj activo de un usuario.

\begin{lstlisting}[language=SQL]
BEGIN;

-- 0. Serializar los cambios de reloj del mismo usuario
SELECT 1 FROM usuarios WHERE AuthorId = :usuario FOR UPDATE;

-- 1. Cerrar la vinculacion activa del usuario (si existe)
UPDATE vinculaciones
SET hasta = now()
WHERE AuthorId = :usuario
  AND hasta IS NULL;

-- 2. Crear la nueva; falla si el reloj ya esta activo en otro
--    usuario (indice unico parcial sobre device_id)
INSERT INTO vinculaciones
    (device_id, AuthorId, desde, hasta, zona_horaria)
VALUES
    (:device_id, :usuario, now(), NULL, :zona_horaria);

COMMIT;
-- Ante cualquier error: ROLLBACK; queda la vinculacion previa
\end{lstlisting}

\textbf{Clasificación: OLTP.} Escribe dos filas (una actualizada y
una insertada). Es infrecuente por usuario, pero concurrente entre
usuarios; latencia objetivo menor a 200 ms.
\textbf{Sin transacción:} si el \codigo{INSERT} fallara tras el
\codigo{UPDATE}, el usuario quedaría sin reloj activo y sus
mediciones no se asignarían a nadie. Si se ejecutaran en orden
inverso sin aislamiento, habría un intervalo con dos relojes activos
y los pasos de ese día se sumarían dos veces.
\textbf{Qué la acelera:} índices únicos parciales
\codigo{(AuthorId) WHERE hasta IS NULL} y
\codigo{(device\_id) WHERE hasta IS NULL}, que además imponen la regla.
\textbf{Concurrencia.} Hay dos mecanismos con funciones distintas.
El \emph{bloqueo} de la fila del usuario (\codigo{SELECT ... FOR UPDATE})
serializa los cambios de un mismo usuario: si llegan dos a la vez, el
segundo espera a que el primero confirme y luego se ejecuta completo,
cerrando la vinculación que el primero acaba de crear e insertando la
suya. Ambos terminan bien y queda activo el último reloj. Los
\emph{índices únicos parciales} garantizan la regla en cualquier caso:
nunca puede haber dos vinculaciones activas del mismo usuario ni del
mismo reloj. Si el reloj nuevo ya está activo en otro usuario, el
\codigo{INSERT} viola el índice sobre \codigo{device\_id}; la
transacción se revierte y la app lo informa.

\subsection{T2. Registrar una comida}

\textbf{Usuario:} la aplicación.

\begin{lstlisting}[language=SQL]
INSERT INTO comidas (comida_id, AuthorId, RecipeId,
                     fecha_hora, registrado_en, porciones)
VALUES (:comida_id, :usuario, :receta,
        :fecha_hora,   -- cuando se preparo (lo indica el usuario)
        now(),         -- cuando se registro en la app
        :porciones);
-- Rechazada por la base si la receta no existe (FK)
-- o si porciones <= 0 (CHECK)
\end{lstlisting}

\textbf{Clasificación: OLTP.} Escribe una fila; frecuente (varias
comidas por usuario al día) y concurrente; menos de 200 ms. Al ser una
sola sentencia, ya es atómica. \textbf{Qué la acelera:} PK y
verificación de FK por índice sobre \codigo{recetas(RecipeId)} y
\codigo{usuarios(AuthorId)}.

\subsection{T3. Reloj activo de un usuario}

\textbf{Usuario:} la aplicación, al abrirse y al sincronizar.

\begin{lstlisting}[language=SQL]
SELECT v.device_id, r.modelo, v.zona_horaria
FROM vinculaciones v
JOIN relojes r ON r.device_id = v.device_id
WHERE v.AuthorId = :usuario
  AND v.hasta IS NULL;
\end{lstlisting}

\textbf{Clasificación: OLTP.} Lee una fila; varias veces por sesión
de uso; alta concurrencia; menos de 100 ms. \textbf{Qué la acelera:}
índice \codigo{(AuthorId, hasta)}.

\subsection{T4. Sesiones recientes de un reloj}

\textbf{Usuario:} la aplicación (historial de actividades).

\begin{lstlisting}[language=SQL]
SELECT session_id, tipo, inicio, fin, distancia_km
FROM sesiones
WHERE device_id = :device_id
ORDER BY inicio DESC
LIMIT 10;
\end{lstlisting}

\textbf{Clasificación: OLTP.} Lee diez filas; frecuente; alta
concurrencia; menos de 100 ms. \textbf{Qué la acelera:} índice
compuesto \codigo{(device\_id, inicio)}, que entrega las filas ya
ordenadas sin recorrer la tabla.

\subsection{A1. Actividad histórica por usuario}

\textbf{Usuario:} analista, y el proceso de entrenamiento para el
modelo de referencia.

\begin{lstlisting}[language=SQL]
SELECT usuario_key,
       AVG(pasos_totales) AS promedio_pasos,
       SUM(pasos_totales) AS pasos_periodo,
       COUNT(*)           AS dias_con_datos
FROM fact_actividad_diaria
WHERE fecha_key BETWEEN :inicio AND :fin
  AND minutos_con_lectura >= :cobertura_minima
GROUP BY usuario_key;
\end{lstlisting}

\textbf{Clasificación: OLAP.} Lee todas las filas usuario-día del
período (unas 18 mil con los datos actuales) y solo cuatro columnas.
Se ejecuta a diario o bajo demanda, con baja concurrencia, y tolera
segundos. \textbf{Qué la acelera:} formato columnar y poda de
particiones por fecha.

\subsection{A2. Actividad y energía de las recetas preparadas}

\textbf{Usuario:} analista. Primero se agregan las comidas por
usuario y día, para no multiplicar las filas de actividad.

\begin{lstlisting}[language=SQL]
WITH comidas_dia AS (
    SELECT usuario_key, fecha_key,
           COUNT(*)         AS numero_comidas,
           SUM(energia_kcal) AS energia_preparada
           -- energia_kcal = porciones * Calories (por porcion)
    FROM fact_comida
    GROUP BY usuario_key, fecha_key
)
SELECT d.fin_de_semana,
       AVG(a.pasos_totales)                  AS pasos_promedio,
       AVG(COALESCE(c.energia_preparada, 0)) AS energia_promedio,
       AVG(COALESCE(c.numero_comidas, 0))    AS comidas_promedio
FROM fact_actividad_diaria a
JOIN dim_fecha d    ON d.fecha_key = a.fecha_key
LEFT JOIN comidas_dia c
       ON c.usuario_key = a.usuario_key
      AND c.fecha_key   = a.fecha_key
GROUP BY d.fin_de_semana;
\end{lstlisting}

\textbf{Clasificación: OLAP.} Combina dos tablas de hechos y una
dimensión y agrega todo el historial (unas 22 mil comidas y 18 mil
usuario-días; millones a escala). Es exploratoria, con baja
concurrencia, y tolera segundos. \textbf{Qué la acelera:} formato
columnar y materializar \codigo{comidas\_dia} como tabla resumen;
A3 usa su versión al corte.

\subsection{A3. Dataset para predecir pasos (responde el objetivo)}

\textbf{Usuario:} proceso de entrenamiento y de predicción diaria.

La consulta usa dos versiones de la tabla de hechos, ambas construidas
en la capa analítica:
\begin{itemize}
    \item \codigo{fact\_actividad\_al\_corte}: para cada fecha $d$, se
    recalcula solo con las lecturas cuya \codigo{ts\_recepcion} es
    anterior al corte (00:00 local de $d+1$). De aquí salen las
    \emph{variables}.
    \item \codigo{fact\_actividad\_diaria}: la versión final, con todas
    las lecturas recibidas. De aquí sale el \emph{objetivo}.
\end{itemize}
Las comidas siguen la misma regla: \codigo{comidas\_dia\_al\_corte}
agrega solo las comidas cuya \emph{hora de registro} en la app es
anterior al corte, aunque se hayan preparado antes. Así, una cena
registrada al día siguiente no entra en la predicción de ese día. Los
datos entregados no traen la hora de registro (\codigo{comidas.csv}
solo tiene \codigo{fecha\_hora} de preparación), por lo que en el
historial se usa \codigo{fecha\_hora} como aproximación. Esa
limitación queda registrada; en producción, la base operacional guarda
la hora de registro.
Además, se genera un \textbf{calendario continuo}: un día por usuario
dentro de su período de vinculación. Un día sin lecturas queda con
pasos nulos y cobertura cero, \emph{no} con cero pasos. Así la ventana
de 7 filas equivale a 7 días calendario, y el objetivo es el día
calendario siguiente, no la siguiente fila con datos.

\begin{lstlisting}[language=SQL]
WITH calendario AS (       -- un dia por usuario dentro de su vinculacion
    SELECT v.usuario_key, f.fecha_key
    FROM periodos_vinculacion v
    JOIN dim_fecha f ON f.fecha BETWEEN v.desde_local AND v.hasta_local
),
base AS (
    SELECT k.usuario_key, k.fecha_key,
           a.pasos_totales,                      -- NULL si no hubo lecturas
           COALESCE(a.minutos_con_lectura, 0) AS cobertura_min,
           COALESCE(a.cantidad_sesiones, 0)   AS cantidad_sesiones,
           COALESCE(c.energia_preparada, 0)   AS energia_preparada
    FROM calendario k
    LEFT JOIN fact_actividad_al_corte a
           ON a.usuario_key = k.usuario_key AND a.fecha_key = k.fecha_key
    LEFT JOIN comidas_dia_al_corte c        -- como A2, con registro < corte
           ON c.usuario_key = k.usuario_key AND c.fecha_key = k.fecha_key
),
variables AS (
    SELECT *,
        AVG(pasos_totales)     OVER w7 AS promedio_pasos_7d,  -- ignora NULL
        SUM(cobertura_min)     OVER w7 AS cobertura_7d,
        SUM(cantidad_sesiones) OVER w7 AS sesiones_7d,
        SUM(energia_preparada) OVER w7 AS energia_preparada_7d
    FROM base
    WINDOW w7 AS (PARTITION BY usuario_key ORDER BY fecha_key
                  ROWS BETWEEN 6 PRECEDING AND CURRENT ROW)
)
SELECT v.*, o.pasos_totales AS pasos_dia_siguiente      -- objetivo
FROM variables v
JOIN fact_actividad_diaria o                             -- version final
  ON o.usuario_key = v.usuario_key
 AND o.fecha_key   = v.fecha_key + 1     -- fecha_key: entero correlativo por dia
WHERE o.minutos_con_lectura >= :cobertura_minima;
\end{lstlisting}

\textbf{Clasificación: OLAP.} Es un batch que ordena el historial de
cada usuario y calcula ventanas. Lee todas las filas usuario-día,
corre a diario o en cada entrenamiento y tolera minutos.
\textbf{Qué la acelera:} partición por fecha, lectura columnar y
materialización de las variables.

\textbf{Días sin cobertura.} Los días cuyo objetivo no alcanza la
cobertura mínima se \emph{excluyen} del entrenamiento: no se sabe
cuántos pasos dio el usuario y tratarlos como cero sesgaría el modelo.
El umbral (minutos con lectura) se fijará en la E2 a partir de la
distribución observada. Los días sin cobertura sí se mantienen en el
calendario, de modo que las ventanas siguen siendo de 7 días.

\textbf{Control de fuga.} Las variables usan solo días hasta $d$ y el
objetivo es el día $d+1$. Las variables salen de
\codigo{fact\_actividad\_al\_corte} y \codigo{comidas\_dia\_al\_corte},
de modo que el modelo ve el día
$d$ incompleto, igual que en producción, donde el 12\,\% de las
lecturas llega después. Los días anteriores de la ventana también se
toman en su versión al corte de cada uno. Eso es conservador: usan
menos datos de los que habría al predecir, pero nunca datos futuros.

\subsection{Consulta fronteriza: promedio reciente en la app}

\begin{lstlisting}[language=SQL]
SELECT promedio_pasos_7d, pasos_previstos_manana
FROM resumen_actividad_usuario     -- publicada por el batch diario
WHERE usuario_key = :usuario;
\end{lstlisting}

El cálculo es analítico (una ventana sobre el historial), pero la
consulta es operacional: la hace cada usuario al abrir la app, con
alta concurrencia y respuesta en menos de 2 s. Se clasifica como OLTP
de lectura sobre un resultado OLAP precalculado. Se ejecuta fuera de
la base operacional, en la tabla resumen o caché, para no competir
con T1--T4.

\section{Normalización de mediciones}\label{an:norm}

\begin{tabularx}{\linewidth}{@{}p{4.2cm}p{5.2cm}X@{}}
\toprule
\textbf{Campo común} & \textbf{Pulso S} & \textbf{FitBand 2}\\
\midrule
device\_id & device\_id & dev\\
ts\_medicion (UTC) & ts\_medicion (ISO 8601) & timestamp (epoch ms)\\
ts\_recepcion (UTC) & ts\_recepcion & received\_at\\
frecuencia\_cardiaca & frecuencia\_cardiaca & fc\\
pasos & pasos & steps\\
calorias & calorias & kcal\\
\bottomrule
\end{tabularx}

\subsection*{Reglas de calidad propuestas (se medirán en la E2)}

\begin{enumerate}
    \item Renombrar los campos al esquema común y convertir las fechas
    a \codigo{timestamp} UTC.
    \item Deduplicar por la \emph{clave lógica}
    (\codigo{device\_id}, \codigo{ts\_medicion}), conservando la
    recepción más temprana. Elimina 345.599 filas. Por comparación,
    deduplicar por \emph{fila completa} (todas las columnas, incluida
    \codigo{ts\_recepcion}) solo eliminaría 45.584: los 300.015
    restantes son reenvíos que llegan con otra hora de recepción.
    En los datos entregados ninguna clave repetida trae valores
    distintos de pasos, calorías o frecuencia cardíaca. Si apareciera
    alguna, se conservaría la recepción más temprana y el conflicto
    se registraría como regla de calidad.
    \item Rechazar pasos o calorías negativos (hoy no hay casos).
    \item Aceptar frecuencia cardíaca nula y reportar la tasa
    (2,3\,\%).
    \item Asignar cada lectura a la vinculación vigente y derivar la
    fecha local con su \codigo{zona\_horaria} (cuatro zonas de EE.\,UU.).
    \item Descartar comidas con \codigo{RecipeId} inexistente (48) y
    sesiones duplicadas (48).
    \item Reprocesar, con sobrescritura de partición, las fechas
    locales afectadas por lotes tardíos.
\end{enumerate}

Pulso S mide cada minuto y FitBand 2 cada cinco; durante una sesión,
ambos miden cada 10 segundos. La cantidad de lecturas no equivale a
la cantidad de pasos: cada lectura trae el incremento desde la
anterior.

\section{Volumen, estimaciones y supuestos}

\begin{tabularx}{\linewidth}{@{}p{4.2cm}X@{}}
\toprule
\textbf{Magnitud} & \textbf{Valor}\\
\midrule
Recetas Food.com & 522.517 (179 MB Parquet).\\
Reseñas Food.com & 1.401.982 de 271.907 usuarios (174 MB Parquet).\\
Relojes & 315 (198 Pulso S y 117 FitBand 2), 300 usuarios,
15 desvinculados.\\
Período & 61 días de recepción, del 16-dic-2007 al 14-feb-2008.\\
Archivos de mediciones & 1.449 archivos horarios .jsonl.gz (115 MB).\\
Mediciones & 15.540.132 filas (15.194.533 tras deduplicar).\\
Sesiones y comidas & 4.840 y 22.106 filas.\\
Crecimiento con 315 relojes & unas 255 mil lecturas al día;
1,9 MB al día y 57 MB al mes, comprimidos.\\
Atraso de recepción & Mediana 1,2 h; percentil 90 de 6,5 h; percentil 99 de 32 h;
máximo 105 h; 12\,\% llega otro día.\\
\codigo{Calories} (Food.com) & Por porción. Si fuera por receta, crecería con
\codigo{RecipeServings}, y no lo hace: la mediana es 227 kcal con 1 porción,
349 con 4 y 156 con 24 (correlación $-0{,}09$).\\
\bottomrule
\end{tabularx}

\textbf{Procedimiento.} Todas las cifras se obtuvieron con consultas
DuckDB sobre los archivos entregados. El script
\codigo{entregas/E1/anexos/inventario\_datos.py} del repositorio las
reproduce: conteo de archivos y filas, unificación de los dos formatos
de reloj, duplicados por clave y por fila completa, atrasos, nulos,
referencias inexistentes y la verificación de \codigo{Calories}.

Para un reloj con $h$ horas al día fuera de sesión y $s$ horas en
sesión, el máximo teórico de lecturas es:
\[
m_{\mathrm{PulsoS}} \approx 60h + 360s,
\qquad
m_{\mathrm{FitBand2}} \approx 12h + 360s,
\qquad h + s \le 24.
\]
El promedio observado (unas 810 lecturas por reloj al día) queda bajo
el máximo de Pulso S (1.440), lo que indica horas sin el reloj puesto.
Por eso la cobertura se incluye como medida.

\clearpage
\section{Registro inicial de decisiones}

{\renewcommand{\arraystretch}{1.25}
\begin{longtable}{@{}>{\raggedright\arraybackslash}p{3.1cm}>{\raggedright\arraybackslash}p{6.1cm}>{\raggedright\arraybackslash}p{8.7cm}@{}}
\toprule
\textbf{Decisión} & \textbf{Alternativa elegida} &
\textbf{Justificación}\\
\midrule
\endhead
Objetivo & Predecir los pasos del día siguiente &
Variable observable; apoya el recordatorio matinal.\\
Operacional & Relacional, con vinculación separada del reloj &
Permite el historial de cambios de reloj y la regla de un activo.\\
Mediciones & Serie temporal fuera de la base OLTP &
15,5 M de filas en 61 días; se leen por reloj y rango de tiempo.\\
Analítico & Esquema en estrella con dos hechos &
Evita mezclar las granularidades usuario-día y comida.\\
Crudos & Archivos originales en objetos &
Trazabilidad y reprocesamiento.\\
Limpios & Parquet particionado por fecha local &
Compresión, lectura columnar y poda por fecha.\\
Arquitectura & Lakehouse &
Separa las capas cruda, limpia y analítica a bajo costo.\\
Escala & Procesamiento local &
115 MB comprimidos; no hay evidencia que justifique distribuir.\\
Deduplicación & Clave (device\_id, ts\_medicion) &
300.015 duplicados llegan con otra hora de recepción.\\
Temporalidad & UTC, más fecha local por vinculación &
Asignar los pasos al día correcto del usuario.\\
Validación ML & Corte por ts\_recepcion y separación temporal &
Evita la fuga por lecturas que llegan tarde.\\
Días sin lecturas & Calendario continuo; objetivos sin cobertura
mínima excluidos &
No confundir un día sin reloj con un día de 0 pasos; ventanas de
7 días calendario.\\
\bottomrule
\end{longtable}
}

\section{Fuentes y referencias}

\begin{enumerate}
    \item INF3862. \emph{Una plataforma de datos para recetas y
    wearables}. Enunciado del proyecto, entregas E1--E3.
    \item INF3862. \emph{Clase 00}. Preguntas guía para el diseño de
    plataformas.
    \item INF3862. \emph{Clase 01: Almacenamiento de Datos}. OLTP, OLAP,
    modelos de datos y arquitecturas analíticas.
    \item Equipo docente. \emph{LEEME.md y diccionario de datos del
    wearable}.
    \item Food.com -- Recipes and Reviews. Kaggle, licencia CC0.\\
    {\small\url{https://www.kaggle.com/datasets/irkaal/foodcom-recipes-and-reviews}}
\end{enumerate}

\end{document}
